\section{A sharper canonical-marginal theorem in every dimension}\label{sec:highdim}
\begin{theorem}[Canonical spectrum and physical Bob spectrum]\label{thm:canonical-spectra}
Let $\rho_B>0$. Let $\alpha_i$ be the decreasing eigenvalues of $\sigma_A$ for the canonical state \eqref{eq:whiten}, and let $\beta_i>0$ be the decreasing eigenvalues of the physical marginal $\rho_B$. Then
\begin{equation}\label{eq:canonical-spectra}
 \boxed{\quad f_d(\rho)\le
 \left(\sum_i\sqrt{\beta_i\min(\alpha_i,1/d)}\right)^2.\quad}
\end{equation}
This bound is sharp for every probability spectrum $\alpha$ and every strictly positive probability spectrum $\beta$.
\end{theorem}
\begin{proof}
Write $\rho=(\Id\otimes K)\sigma(\Id\otimes K^\dagger)$ with $K=\sqrt{d\rho_B}$. For any maximally entangled $e$, put $x=(\Id\otimes K^\dagger)e$. Its squared norm is $\tr(KK^\dagger)/d=1$, and its reduced Bob spectrum is $\beta$. Apply Theorem~\ref{thm:weighted} to $\sigma$, whose marginals have spectra $\alpha$ and $(1/d,\ldots,1/d)$, and to this $x$.

For sharpness, take the construction \eqref{eq:coherent-min} for marginal spectra $\alpha$ and $(1/d,\ldots,1/d)$, and filter with
\[
 K=\diag\bigl(\sqrt{d\beta_1},\ldots,\sqrt{d\beta_d}\bigr).
\]
The resulting state is normalized, has physical Bob spectrum $\beta$, and whitening returns the chosen canonical state. The Bell target $d^{-1/2}\sum_i\ket{ii}$ has expectation equal to the right side of \eqref{eq:canonical-spectra}.\end{proof}

Let $m_i=\min(\alpha_i,1/d)$, $L=\sum_i m_i>0$, and $\nu_i=m_i/L$. Then \eqref{eq:canonical-spectra} reads
\begin{equation}\label{eq:affinity}
 f_d(\rho)\le L\left(\sum_i\sqrt{\beta_i\nu_i}\right)^2\le L.
\end{equation}
The last upper bound is the old unweighted marginal-spectrum result of Klyachko \cite{klyachko}. The retained overlap factor is the additional information here. Equality in its Cauchy--Schwarz step requires $\beta=\nu$. If every $m_i>0$, this is an admissible full-rank Bob spectrum and is attained by the construction. If some $m_i=0$, $L$ is a supremum over positive Bob spectra, approached as they tend to $\nu$; it cannot be attained within that full-rank stratum.

\begin{corollary}[Marginal-spectrum stability]
If $L-f_d(\rho)\le\varepsilon$, then
\begin{equation}\label{eq:highdim-stability}
 \|\beta-\nu\|_1\le2\sqrt{\varepsilon/L}.
\end{equation}
\end{corollary}
\begin{proof}
Write $A=\sum_i\sqrt{\beta_i\nu_i}$. From \eqref{eq:affinity}, $1-A^2\le\varepsilon/L$. Also
\[
 \sum_i|\beta_i-\nu_i|
 \le\sqrt{\sum_i(\sqrt{\beta_i}-\sqrt{\nu_i})^2}
      \sqrt{\sum_i(\sqrt{\beta_i}+\sqrt{\nu_i})^2}
 =2\sqrt{1-A^2}.
\]
This proves the assertion. It is a statement about marginal eigenvalues, not trace-distance closeness of the entire bipartite state.\end{proof}

